% Chapter: review of existing internal tools
\chapter{The tools already in the repository}
\label{sec:tools}

An author working in this workspace already has a prose policy, three skills, a
pattern list, two grounding tools, and a way to distribute them. Those pieces
can help, but they do not yet form a product. We need to know which ones help an
author in practice, which have evidence behind them, and where each one stops.
The rules are mature and largely consistent; executable measurement, behavioral
evidence, and one product-owned home are still missing.

Figure~\ref{fig:tool-gap} explains why consolidation alone is not the product.
The missing value lies in the handoffs between capable tools and in the reader
outcome at the end.

\hvFigureToolGap

\section{The prose policy}
\label{sec:tool-policy}

Every agent receives roughly 150 lines of the global policy under the heading
``Reader-Facing Scientific Prose.''\footnote{\path{claudecodex/policies/global-scientific-coding-agent-policy.md}. The companion section ``Plain Scientific Language And Non-Evasion'' supplies the verdict vocabulary discussed below.}
Its load-bearing device is a precedence order. When rules conflict, an agent
resolves them in this sequence: mathematical and scientific correctness
first, then source and evidentiary fidelity, then domain-appropriate meaning,
then reader comprehension and natural expression, and only then template
regularity. The order is what keeps a pattern list from doing damage. A rule
like ``remove em dashes'' sits at the lowest level, so it can never justify
deleting a numeric range from a theorem statement; a readability preference
can never justify dropping an assumption.

Three rules came directly from the failure record.

The first keeps internal workflow language out of ordinary exposition. Words
such as ``artifact,'' ``pipeline,'' ``gate,'' and ``handoff'' remain available
when the text is genuinely about software or project operations. They are not
used as metaphors for a theorem, an argument, or a paragraph. This rule turns
the administrative-register finding in Table~\ref{tab:bgs-taxonomy} into a
choice an editor can make sentence by sentence.

The second replaces defensive disclaimers with reasons. The policy prefers
``We assume log utility in this example because it keeps the first-order
condition simple'' to ``We are not claiming that households literally have
logarithmic utility.'' The first sentence states the modelling choice and its
purpose. The second introduces a mistaken interpretation merely to deny it.
This is the repair observed in Section~\ref{sec:bgs-audit}, expressed as a
general rule.

The third reserves acceptance for a reader. Author self-review, a model review,
a checklist, a style score, and a clean compilation can each provide evidence
about one part of a document. None can certify a human voice. The vocabulary
beside that rule is deliberately direct: a statement may be correct, wrong for
the stated target, unsupported, not checked, or heuristic. A polite synonym
must not blur which verdict applies.

Merits: the policy is complete in the sense that every failure class in
Chapter~\ref{sec:record} has a corresponding clause, and the precedence order
resolves the conflicts that sank naive rule-following elsewhere. Problems:
it is prose all the way down. Nothing in it executes, so compliance is
whatever the drafting agent happens to retain from 618 lines of context, and
the one thing the policy cannot do is notice that it is being ignored. The
ZLB rescue's third pass began with the admission that the policy ``had not
been consulted.'' A policy that must be remembered is a policy that will
sometimes not be.

\section{The narrative skill}
\label{sec:tool-linear}

When an agent uses \path{write-linear-scholarly-narrative}, it gets the
document-architecture skill,
distilled from the SMEwallet Draft~11 method.\footnote{\path{claudecodex/skills/write-linear-scholarly-narrative/SKILL.md}, with a 259-line companion, \path{references/teaching-contract.md}, defining the comprehension standard.}
Its sequence begins with the reader rather than the source files. The writer
states who will read the unit, what that person already knows, and what they
should be able to explain afterward without hunting through another document.
For this proposal, that means a former professor acting as an executive, not an
engineer who already knows the repository.

The writer then chooses the smallest case or question that can carry the
argument. Chapter order follows explanatory dependence: a concept appears
after the reader has encountered the problem that requires it, even if the
source material was originally stored in another order. Each bounded unit
opens one question and introduces the smallest idea needed to answer it.

Drafting follows a teaching sequence. Return to a known example, show the new
difficulty, and introduce notation only when it resolves that difficulty. An
equation is followed by its interpretation. The transition to the next unit
comes from a question the present unit cannot yet answer, not from a stock
phrase announcing that the document will move on.

Finally, a person reads the unit and tries to reconstruct the chain. They
should be able to explain why the unit began, what each turn contributed, what
the equation was for, and why the next question follows. Failure at that stage
is evidence about the exposition, not a request for the reader to study harder.

Its most valuable property is self-awareness about
Section~\ref{sec:lessons}'s fourth lesson. Nearly every rule carries an
anti-template clause; the drafting sequence is offered ``as a menu,'' with
the instruction not to fill slots mechanically and not to add a caveat or
transition merely because the template has a place for one. The
working-memory budget (about four live conceptual items) is labeled a
diagnostic threshold, not a law. Whoever wrote it had watched a floor become
a ceiling.

Evidence of effect is real but bounded: Draft~11's units passed their PI
gates one by one, which no earlier method achieved. Problems: the skill has
no examples inside it, so an agent meets its abstractions (``narrative
spine,'' ``reactivate'') without a single worked instance; it has no tooling,
so storyboards and contracts live or die by agent diligence; and its
acceptance step needs a human per unit, which priced it out of the BGS lane's
volume. Nothing tests whether an agent following it writes measurably better
than an agent that has never seen it, a gap Chapter~\ref{sec:proposal}'s
third work package exists to close.

\section{The sentence-level skill}
\label{sec:tool-natural}

The companion \path{write-natural-analytical-prose} works at sentence and
paragraph scale.\footnote{\path{claudecodex/skills/write-natural-analytical-prose/SKILL.md}, with before/after contrasts in \path{references/prose-patterns.md}.}
Before any rewrite, the agent
freezes a meaning ledger: the facts, attributed claims, deductions, live
uncertainties, author judgments, and exact numbers that must survive. After
the rewrite, it diffs the result against the ledger, and a smoother sentence
that changed a claim counts as a failure, not an improvement. Between those
two steps sit the craft rules: give each sentence one epistemic job, with a
calibrated verb table in which ``shows'' is not a synonym for ``suggests''
and ``reproduces'' is not a synonym for ``validates''; put concrete actors in
subject position (``the regime test rejects the authors' path,'' not ``the
evidence indicates''); state the finding before its qualifications; and spend
each qualification once, at the point of likely misunderstanding, rather than
re-attaching it to every positive sentence.

This is the skill that operationalizes the epistemic-blur and abstract-actor
rows of Table~\ref{tab:bgs-taxonomy}, and its meaning-ledger loop is the
sentence-scale version of the protected-baseline invariant. Its problems
mirror the narrative skill's: no automation, so the ledger diff is manual;
verification by ``read it aloud''; and hard prohibitions that depend entirely
on the drafting agent noticing its own violations. The skill is also, like
everything in this stack, tested only by phrase-lock (see
Section~\ref{sec:tool-claudecodex}): the test suite asserts that certain
sentences still appear in the skill file, not that the skill changes any
output.

\section{The humanizer pattern list, as installed}
\label{sec:tool-humanizer}

On 20~August the workspace vendored \texttt{blader/humanizer} v2.11.2 into
the policy tree, as
\path{humanizer-ai-writing-patterns.md}.\footnote{\path{claudecodex/policies/humanizer-ai-writing-patterns.md}. Section~\ref{sec:oss-humanizer} reviews the upstream project; this subsection reviews our installation of it.}
The list itself is 35 numbered AI-writing patterns derived from Wikipedia's
``Signs of AI writing'' catalogue, each with words to watch and a
before-and-after pair; Section~\ref{sec:oss-humanizer} walks through its
content. What matters here is the installation decision. The file carries a
locally written precedence header that demotes the whole list to ``a
DIAGNOSTIC pattern list, subordinate to
global-scientific-coding-agent-policy.md,'' spells out the order of
authority, and singles out the one rule that would otherwise do damage: the
flat dash ban of its pattern~14 is to be applied ``proportionately rather
than as a ban,'' and no pattern may be satisfied ``by deleting mathematics,
assumptions, or qualifications.''

That header exists because of a live incident. During the ZLB rescue the
list was consulted without a precedence frame, and the case study records
the realization that without one, ``a pattern list becomes a new template and
you trade one uniformity for another.'' The header is currently the only
policy-precedence mechanism in the workspace, and it is a comment in a
markdown file. It governs an agent exactly as far as the agent reads and
honors it. Turning this from a comment into configuration, per-format rule
downgrades in the diagnostic layer itself, is one of the cleaner design
requirements to come out of this review
(Chapter~\ref{sec:design}, under configurable policy precedence).

\section{The repair skill and the only working scripts}
\label{sec:tool-repair}

MathDevMCP ships its own scholarly-writing skill,
\path{repair-scholarly-readability}, and it is the only piece of the stack
with executable parts.\footnote{\path{MathDevMCP/skills/repair-scholarly-readability/SKILL.md}, with three reference rubrics and two scripts under the same directory.}
Where the narrative skill teaches drafting, this one teaches repair of an
existing LaTeX document, and its structure encodes two ideas the rest of the
stack lacks. The first is four separate records: reader comprehension,
mathematical integrity, source fidelity, and typography, with the explicit
rule that ``evidence in one ledger cannot certify another.'' A clean compile
is typography evidence only; a passed equation audit says nothing about
motivation. The second is bounded repair: one mechanism chapter or 6--12
rendered pages at a time, against a preserved baseline, ending in a test with a
reader who has no project history, prespecified questions, and a four-record
decision table.

Both scripts matter beyond their modest size.
\path{scripts/scholarly_surface_audit.py} emits JSON and markdown surface
diagnostics for a .tex/.pdf pair, and \path{scripts/render_review_pages.sh}
renders page images for visual inspection, institutionalizing the ZLB
rescue's discovery that text extraction lies about layout. These are the
embryo of the \texttt{hv capture} command in Chapter~\ref{sec:proposal}.

Evidence of effect: the skill's bounded method took the opening of the BGS
industry-DSGE dossier through a repair that passed five prespecified
reader-without-project-history questions, the strongest automated-plus-human result in the BGS
lane. The same result file records the limit: later chapters ``remain dry,''
and the project's audit concluded the skill ``generalizes locally, not across
a corpus-sized argument,'' which is lesson five of
Section~\ref{sec:lessons} observed in the field. The other problem is
organizational: this skill overlaps the claudecodex narrative skill in scope
while living in a different repository with different reference rubrics.
Scholarly readability currently has two homes and no owner.

\section{ResearchAssistant}
\label{sec:tool-ra}

When a citation needs checking, ResearchAssistant supplies the source-grounding
half of the discipline. It is a local-first tool that ingests papers (preferring arXiv LaTeX source over PDF,
because the LaTeX parses reliably and the PDF does not), extracts sections,
equations, labels, and bibliographies, tracks provenance and review status,
and exposes a read-only MCP surface plus a
CLI.\footnote{\path{research-assistant/}; the checked-in MCP configuration targets Claude Code and VS Code. Python 3.11 only.}
For this survey its relevant property is precedent: its \texttt{ra survey}
subcommands are the one place where a piece of workspace writing policy
became software. The literature-audit policy demands six ledgers (source
support, citation metadata, backward and forward snowballing, claim support,
omitted-paper risk); \texttt{ra survey central-papers} writes exactly those
ledgers and returns per-paper dispositions.

Its behavior under failure is the part worth copying. In the ZLB campaign
the tool's public-discovery bootstrap was unavailable, and instead of
degrading silently it closed the mission as
\texttt{terminal\_blocked\_bootstrap\_unavailable}, leaving the ledgers to be
maintained by hand. A tool that reports its own incapacity honestly is rarer
than it should be. Its limits for \hv{} are equally clear: it grounds
sources, not prose; its equation extraction from PDFs is explicitly marked
unreliable (arXiv LaTeX is the reliable path); and its discovery pipeline
depends on external metadata services that were down when most needed.

\section{The mathematics and modeling tools}
\label{sec:tool-math}

Two more tools border the writing problem without being writing tools.
MathDevMCP's mathematical layer provides SymPy-backed equality checking,
derivation audits, and an experimental whole-document rigor audit over
LaTeX. Its one recorded contribution to a manuscript is instructive about
scale: the ZLB audit produced per-equation issue ledgers and exactly one
applied exposition patch (an invertibility condition stated before a
displayed inverse), tagged
``candidate\_exposition\_patch\_not\_certificate.''\footnote{\path{BayesFilter/docs/surveys/zlb_discontinuous_hmc/mathdevmcp_audit_20260819.md}.}
That is the correct division of labor, mathematics tools nominating exposition
fixes without certifying prose, and \hv{} should preserve it rather than
absorb it. DynareMCP, the macro-model execution server, contributes structure
inventories that were used to falsify inflated rewrite claims (counting
inherited words and relocated blocks), a niche but genuinely useful
falsification tool for ``we rewrote it'' assertions.

\section{Distribution and the limits of the current tests}
\label{sec:tool-claudecodex}

claudecodex is the delivery mechanism. Its installer splices the global
policy into marked blocks in each agent's configuration files idempotently,
merges approval rules, and copies the skills into the Codex skill
tree.\footnote{\path{claudecodex/install_global_agent_policy.py}; templates and review-gate scripts under \path{claudecodex/docs/templates} and \path{claudecodex/scripts}.}
Around it sit the cross-agent review plumbing (bounded review packets, probe
ladders, a gate script that parses ``VERDICT: AGREE or REVISE'') and a
five-file proposal template system with seven evidence ledgers. The
machinery works and is in daily use.

Its test suite deserves a hard look because of what it teaches about the
stack's blind spot. There are real behavioral tests for the review gate (a
fake worker exercises the probe-and-fallback paths) and for the installer.
For the writing skills there are only phrase locks:
\path{test_natural_analytical_prose_skill.py} asserts that certain sentences
still occur in the skill file and that banned phrasings do not. Phrase locks
prevent silent regression of the doctrine's text. They say nothing about
whether the doctrine changes what an agent writes. As of today, no experiment
anywhere in the workspace has compared documents produced with and without
any of these skills. The doctrine is plausible, consistent, and untested.

Two further operational gaps matter. Installation is
Codex-first; Claude Code agents receive the policy but not a native skill
installation. And the writing doctrine now lives in three places, the
claudecodex policy tree, the claudecodex skills, and the MathDevMCP skill,
with cross-references in prose instead of a shared source of truth.

\section{The gap, stated precisely}
\label{sec:tool-gap}

The last column of Table~\ref{tab:tool-gap} gives the build list. The six
capabilities requested by the ZLB case exist nowhere as a coherent product: a
register linter, rhythm profiler, defensive-register finder, LaTeX-aware
baseline differ, a private provenance map that keeps evidence attached to one
manuscript, and a precedence mechanism stored in configuration rather than a
comment. The doctrine also has no behavioral evidence, and its three copies
need a shared source. Humanvoice joins that list into one testable review path.

\begin{table}[t]
\centering
\small
\begin{tabularx}{\textwidth}{@{}l L L@{}}
\toprule
Piece & What it reliably provides & What it cannot do \\
\midrule
Prose policy & complete doctrine; precedence order & execute, or notice being ignored \\
Narrative skill & Draft-11 method; anti-template clauses & show examples; run without a human per unit \\
Sentence skill & meaning-ledger safety loop; verb calibration & automate the ledger diff \\
Humanizer (installed) & 35-pattern diagnostic with FP guards & bind its own precedence outside prose \\
Repair skill & four-ledger repair; the only scripts & compose beyond 6--12-page units \\
ResearchAssistant & source grounding; ledgers as software & anything about prose \\
MathDevMCP/DynareMCP & equation audits; structure inventories & certify exposition \\
claudecodex & idempotent install; review plumbing & test writing behavior; reach Claude natively \\
\bottomrule
\end{tabularx}
\caption{The existing stack, by what each piece does and does not provide.}
\label{tab:tool-gap}
\end{table}

The inventory is therefore an asset and a warning. It gives the build a set of
tested interfaces, historical fixtures, and language for describing failure;
it does not give the project a validated intervention. The next part places
those local assets beside external evidence. That comparison matters because
a tool can be useful in this repository for reasons that do not travel to
another genre, and a widely used package can still fail the protected-source
contract. The field review begins with what research supports, then explains
how the literature was assembled, before returning to the package choices.
